\chapter{Foreword}
This book is a collection of knowledge, notes, and any learning material I had to use to learn or recall on topics like continuum mechanics, machine learning, nonlinear FEA and the use of open source softwares like libROM and MFEM. The end goal is to have a clean notebook with what I learned or what at I had to recall to implement my end goal: \textbf{\textit{A Framework for Hyper-reduction of Viscoplastic Models applied to high temperature components}}. While I took a lot of my free time (and cut on sleep once in a while), worked hard to test and verify to make the most out of this book it is most likely certain that it contains some error due to either formatting and/or comprehension. I used Large Language Models (LLMs) to answer questions or verify certain parts of the text, but still, I think some errors may have slipped in. It is not a sponsored project, this is something I worked on during my free time because I am genuinly interested by te topics adresses in this book. In no case this book should be taken as is. Each and every cited ressources should be read to fully understand the notes I took. I take no responsibility regarding the exactitude of any equations or any theory related claim. I would be happy to help someone learn, but the learner should help him/her self and take appropriate actions to validate his/her learning. 

 This project started because I wanted to learn and work on something relevant to the current finite element analysis (FEA) / reduced order models (ROM) / artificial intelligence (AI) landscapes. I think that AI's biggest problem related to simulation is data reliance. While technologies like Physics Informed Neural Networks (PINNs) are good specific machine learning techniques applied to simulation, other technologies still rely heavily on data to calibrate models. In this data driven world, I couldn't stand companies selling :\textit{AI results in minutes rather than days!} when in fact, it took most likely days and numerous well thought and well built models to feed the "super fast" AI prediction. But thats when I had the idea to first evalutate the capability of a "non AI" solution to predict results of large scale, highly nonlinear problems. Maybe to even solve complex problem faster and with higher accuracy than an AI approach? Or maybe feed AI models more efficiently?  Then I started reading on the topics and ask Claude (by Anthropic) to act as a supervisor and guide me through this quest to learn and find interestings things.
\\

\noindent
Billy Guy, M.Sc.A